% Biên dịch bằng XeLaTeX hoặc Tectonic.
\documentclass[10pt,a4paper]{article}
\usepackage{fontspec}
\setmainfont{texgyretermes}[Extension=.otf,UprightFont=*-regular,BoldFont=*-bold,ItalicFont=*-italic,BoldItalicFont=*-bolditalic]
\usepackage[margin=1.9cm,headheight=14pt]{geometry}
\usepackage{amsmath,amssymb,booktabs,tabularx,array,ragged2e,enumitem,microtype,xcolor,tikz,fancyhdr}
\usetikzlibrary{arrows.meta,positioning}
\usepackage[hidelinks,unicode]{hyperref}
\definecolor{ink}{HTML}{183A45}
\definecolor{teal}{HTML}{087F7A}
\definecolor{pale}{HTML}{EAF4F2}
\definecolor{gray}{HTML}{56656B}
\newcolumntype{Y}{>{\RaggedRight\arraybackslash}X}
\newcolumntype{P}[1]{>{\RaggedRight\arraybackslash}p{#1}}
\setlength{\parindent}{0pt}\setlength{\parskip}{5pt}
\setlength{\emergencystretch}{2em}
\renewcommand{\arraystretch}{1.17}
\setlist{nosep,leftmargin=1.4em,topsep=4pt}
\pagestyle{fancy}\fancyhf{}
\fancyhead[L]{\small\color{gray}BẢO VỆ RIÊNG TƯ QUỸ ĐẠO}
\fancyhead[R]{\small\color{gray}19/09/2026}
\fancyfoot[C]{\small\thepage}
\renewcommand{\headrulewidth}{0.3pt}
\newcommand{\key}[1]{\par\smallskip\noindent\colorbox{pale}{\parbox{\dimexpr\linewidth-2\fboxsep}{\textbf{#1}}}\par\smallskip}
\newcommand{\up}{\ensuremath{\uparrow}}\newcommand{\down}{\ensuremath{\downarrow}}
\hypersetup{pdftitle={Bảo vệ riêng tư quỹ đạo: khung bài toán và giao thức đánh giá},pdfauthor={}}
\begin{document}
\begin{center}
{\LARGE\bfseries\color{ink} Bảo vệ riêng tư quỹ đạo}\\[5pt]
{\large Khung bài toán, đối chứng và giao thức đánh giá}\\[4pt]
{\small 19/09/2026}
\end{center}
\textbf{Mục tiêu.} Xây dựng cơ chế bảo vệ cho xe đô thị sử dụng dịch vụ tìm địa điểm (POI), đồng thời duy trì chất lượng dịch vụ và chi phí phù hợp. Nghiên cứu đặc tả S1--S10 cho bốn mục tiêu riêng tư; giai đoạn phát triển và so sánh đầu tiên tập trung S1--S3.

\section{Bài toán, kịch bản và dữ liệu}
\subsection{Bốn mục tiêu bảo vệ}
\begin{tabularx}{\linewidth}{@{}P{3.1cm}Y P{3.4cm}@{}}
\toprule\textbf{Mục tiêu}&\textbf{Thông tin cần giữ riêng}&\textbf{Kịch bản}\\\midrule
Vị trí và quỹ đạo & Đang ở đâu, dừng ở đâu, đường đã đi, điểm đầu và điểm cuối. & S1--S3, S9--S10; S8 bổ sung dữ liệu đồng hành.\\
Danh tính & Các phiên có thuộc cùng người hoặc thiết bị không. & S4.\\
Tuyến tiếp theo & Cạnh đường sắp đi và đích chưa tới. & S5--S6.\\
Nội dung truy vấn & Loại dịch vụ và ý định thật của người dùng. & S7.\\\bottomrule
\end{tabularx}

\subsection{Từ mục tiêu đến benchmark}
Đối thủ cơ sở là nhà cung cấp dịch vụ vị trí (LSP): xử lý đúng yêu cầu nhưng lưu và phân tích dữ liệu nhận được. Đối thủ biết thuật toán, bản đồ và tham số công khai. Lịch sử liên phiên, hồ sơ và dữ liệu đồng hành chỉ được cấp theo từng phép thử. Đáp án thật và trạng thái riêng trên thiết bị thuộc phía đánh giá.

\begin{center}
\begin{tikzpicture}[box/.style={draw=teal,fill=pale,rounded corners=2pt,text width=3.25cm,minimum height=1.15cm,align=center,font=\small},>=Latex]
\node[box] (a) {\textbf{4 mục tiêu}\\Thông tin cần giấu};
\node[box,right=0.55cm of a] (b) {\textbf{10 kịch bản}\\Quan sát, phụ trợ, đáp án};
\node[box,right=0.55cm of b] (c) {\textbf{Dataset có nhãn}\\Chia học/chọn/kiểm thử};
\node[box,below=0.7cm of c] (d) {\textbf{Benchmark}\\Đối thủ và bộ đo};
\node[box,left=0.55cm of d] (e) {\textbf{Phát triển S1--S3}\\Tối ưu và đối chiếu};
\node[box,left=0.55cm of e] (f) {\textbf{Mở rộng S4--S10}\\Bổ sung cơ chế phù hợp};
\draw[->] (a)--(b);\draw[->] (b)--(c);\draw[->] (c)--(d);\draw[->] (d)--(e);\draw[->] (e)--(f);
\end{tikzpicture}
\end{center}
\textbf{Hình 1.} Kịch bản quyết định dữ liệu và quyền quan sát trước khi xây dựng các thuật toán tấn công. Survey cung cấp cơ chế và các giới hạn cần đưa vào phép thử.

\subsection{Đặc tả chung cho mỗi phép thử}
Mỗi bản ghi cần xác định: (1) thông tin cần suy ra; (2) phần transcript được xem, gồm thời điểm, tọa độ/tập tọa độ, định danh và nội dung công bố; (3) thông tin phụ trợ; (4) thời điểm suy luận; (5) đáp án và đơn vị chấm. S1--S10 là hệ kịch bản của nghiên cứu, mỗi kịch bản có ba ca A/B/C để kiểm tra điều kiện khác nhau.

\key{Chuẩn hóa nhiệm vụ và quyền quan sát; dùng nhiều đối thủ phù hợp với từng giao diện đầu ra.}
Một phương pháp công bố điểm thay thế không có ``thành viên thật'' để đối thủ chọn. Đối thủ phải suy ra cùng loại đáp án từ transcript thực sự nhận được. Mô hình tấn công được học và chọn trên dữ liệu riêng, sau đó cố định khi kiểm thử.

\newpage
\subsection{Mười kịch bản và quyền quan sát}
\par\small
\begin{tabularx}{\linewidth}{@{}P{0.65cm}P{2.45cm}Y P{3.45cm}@{}}
\toprule\textbf{Mã}&\textbf{Đáp án cần suy ra}&\textbf{Transcript và thông tin phụ trợ}&\textbf{Ba ca kiểm tra}\\\midrule
S1 & Vị trí tại lần gửi & Một bản tin; bản đồ và thống kê công khai từ tập học. & Nhiều hướng đi; một hướng; gần loại POI hiếm.\\
S2 & Vị trí dừng & Nhiều bản tin trong cửa sổ dừng; bản đồ và thống kê nền. & Dừng ngắn; dừng dài; rời đi rồi quay lại.\\
S3 & Đường đã đi & Chuỗi công bố trong cửa sổ; mạng đường và mô hình chuyển động từ tập học. & Qua nhiều đoạn; hành lang ít lựa chọn; gửi thưa.\\
S4 & Cùng người / thiết bị & Hai phiên không chồng lấn; đặc trưng/định danh thực sự công bố. Hồ sơ định danh chỉ cấp nếu phép thử quy định. & Cùng người, cùng máy; cùng người, đổi máy; khác người, chung máy.\\
S5 & Cạnh đường kế tiếp & Chỉ tiền tố trước bước chuyển; tập cạnh hợp lệ trên bản đồ. Đoạn tương lai giữ làm đáp án. & Có lựa chọn rẽ; một lựa chọn; cùng tiền tố, khác bước tiếp.\\
S6 & Đích chưa tới & Tiền tố trước khi tới đích; ca C cấp thêm sáu chuyến lịch sử liên quan đã qua bảo vệ. & Chung tiền tố, khác đích; hai đích gần; đích quen/hiếm.\\
S7 & Loại nhu cầu / ý định & Nội dung và chuỗi truy vấn được công bố; nhãn ý định thật được giữ kín. & Nội dung nguyên văn; trộn sáu loại; cùng truy vấn đầu, khác chuỗi sau.\\
S8 & Vị trí khi có đồng hành & Transcript mục tiêu; bật/tắt dữ liệu người đồng hành và quan hệ theo cấu hình. & Đi cùng rồi tách; chung tuyến; tình cờ ở gần.\\
S9 & Điểm xuất phát bị che & Phần công bố sau đoạn đầu bị ẩn; bản đồ. Metadata chỉ dùng khi được khai báo. & Một lối ra; khác điểm đầu rồi nhập tuyến; lặp điểm xuất phát.\\
S10 & Điểm kết thúc bị che & Phần công bố trước đoạn cuối bị ẩn, sau khi chuyến kết thúc; bản đồ. & Một lối vào; chung tuyến rồi tách đích; lặp điểm kết thúc.\\\bottomrule
\end{tabularx}
\normalsize

\textbf{Phân biệt thời điểm suy luận.} S5--S6 dự đoán tương lai khi chuyến còn diễn ra. S9--S10 suy ngược phần đã xảy ra nhưng bị che. S4 xét liên kết phiên; S7 xét nội dung; S8 xét tác động của quan sát đồng hành.

\subsection{Dataset và phân chia đánh giá}
Dữ liệu dùng SUMO trên mạng đường OSM, gồm quỹ đạo, thời gian, truy vấn POI và nhãn kịch bản. Phiên bản khảo sát có \textbf{12 nhóm tuyến, 264 chuyến hoàn tất và 393 bản ghi}, phủ 30 ca A/B/C. Một chuyến có thể tạo nhiều bản ghi; các bản ghi liên quan phải được kiểm soát khi chia dữ liệu và tổng hợp kết quả.

\begin{itemize}
\item \textbf{Tập học:} ước lượng thống kê vị trí/chuyển tiếp, huấn luyện cơ chế và đối thủ.
\item \textbf{Tập chọn:} chọn tham số, đối thủ và cấu hình theo tiêu chí đã định.
\item \textbf{Tập kiểm thử:} xác nhận sau khi khóa các lựa chọn; không dùng lại để điều chỉnh mô hình.
\end{itemize}
Nhãn danh tính, ý định và quan hệ đồng hành hiện được xây dựng trong mô phỏng; S7 dùng mẫu ý định viết tay. Cần kiểm tra độ phủ và độ cân bằng trước khi mở rộng kết luận sang S4--S10.

\newpage
\subsection{Các hướng nghiên cứu và mức độ giải quyết}
Bảng dưới tổng hợp cơ chế và mức phù hợp với từng kịch bản. Các công trình nền tảng được giữ cùng nghiên cứu gần đây để thể hiện sự phát triển của hướng giải quyết.
\par\small
\begin{tabularx}{\linewidth}{@{}P{0.65cm}P{3.25cm}Y Y@{}}
\toprule\textbf{Ca}&\textbf{Hướng / reference}&\textbf{Cách bảo vệ và ưu điểm}&\textbf{Giới hạn và nhận định}\\\midrule
S1 & DLS/enhanced-DLS (2014) [1]; Geo-I (2013) [2] & Chọn dummy có độ phổ biến tương tự, tăng phân tán; hoặc làm nhiễu vị trí theo bảo đảm phân biệt. & Có cơ sở cho bảo vệ từng lần gửi; cần đánh giá tích lũy khi lặp lại.\\
S2 & ASA (2017) [3]; cache TPP-CCRS (2024) [4] & Cân bằng thống kê dummy/bí danh; cache giảm nhu cầu gửi lại. & Chống thống kê và giảm công bố là hai hướng bổ sung; cache có yêu cầu kiến trúc, độ mới dữ liệu. Chưa đồng nhất phép thử dừng của S2.\\
S3 & RDG (2021) [5]; TransProtect (2024) [6]; semantic correlation và fake queries (2026) [7,8] & Xét chuyển tiếp, ngữ cảnh đường, ngữ nghĩa hoặc chèn truy vấn để làm khó nối chuỗi. & Có nhiều cơ chế và thực nghiệm theo đối thủ riêng; cần so sánh cùng utility và chi phí. Là trọng tâm đối chứng giai đoạn đầu.\\
S4 & Mix Zones (2004) [9]; CPCROK (2025) [10] & Đổi bí danh trong vùng trộn; bổ sung dự đoán và quỹ đạo giả cho mạng xe thưa. & Hiệu quả phụ thuộc trộn và thông tin liên kết. VANET khác LSP; cần chấm người/thiết bị riêng.\\
S5 & Hide-and-Seek (2014) [11]; chống dự đoán, Qiu et al. (2023) [12] & Tính tương quan tương lai; điều chỉnh bảo vệ theo khả năng bị dự đoán. & Có cơ chế chống suy luận theo thời gian; cần đối chiếu đúng tác vụ cạnh kế tiếp và quyền xem tiền tố.\\
S6 & DesPre/CkiDel (2017) [13]; DFPS [14] & Xóa check-in làm lộ đích; hoặc che đích trong vùng nhiều lựa chọn khi tìm chỗ đỗ. & Bảo vệ đích trong từng ứng dụng; chưa đủ chứng minh chống suy đích từ luồng xe trực tuyến của S6.\\
S7 & Chuỗi truy vấn giả (2021) [15]; bảo vệ vị trí--truy vấn cho xe (2024) [16] & Tạo chuỗi che cả tọa độ và nội dung, khai thác tương quan POI. & Đúng mục tiêu; cần kiểm tra ý định dài hạn và overhead, không chỉ từng loại truy vấn.\\
S8 & Phối hợp làm mờ (2017) [17]; DP-FETC (2025) [18] & Phối hợp dữ liệu đồng vị trí; bảo vệ tương quan khi công bố quỹ đạo. & Bảo vệ riêng một người không kiểm soát mọi tiết lộ từ người khác. DP-FETC là hướng công bố dữ liệu, chưa là đối chứng trực tuyến.\\
S9 & Tấn công EPZ và biện pháp giảm rò rỉ (2022) [19] & Che đoạn rời điểm đầu, kết hợp kiểm soát metadata. & Ranh giới che và mạng đường vẫn là manh mối; che đầu đơn thuần chưa đủ trước tấn công đã khảo sát.\\
S10 & EPZ (2022) [19] & Che đoạn tới điểm cuối; hạn chế metadata hỗ trợ suy ngược. & Cần kiểm tra nhiều chuyến và dữ liệu phụ trợ; khác bài toán dự đoán đích S6.\\\bottomrule
\end{tabularx}
\normalsize

\textbf{Hướng mở rộng khảo sát.} S3 đã có nhiều cơ chế; các nhóm còn lại cần đọc sâu và đối chiếu thêm thực nghiệm, mã nguồn và điều kiện triển khai. Mỗi kịch bản ưu tiên một công trình nền tảng, một hướng khác hoặc gần đây và bằng chứng về giới hạn. Các nguồn [4,10,12,14,16,18] mở rộng góc nhìn, chưa mặc định trở thành đối chứng chạy trực tiếp.

\textbf{Bài học thiết kế.} Dummy cần hợp cả không gian và thời gian; nội dung và định danh cần cơ chế riêng; bảo vệ đồng hành đòi hỏi kiểm soát tương quan; che biên cần xét cả metadata. Đây là cơ sở mở rộng phương pháp sau S1--S3.

\newpage
\section{Đối chứng và bộ metrics chung}
\subsection{Năm đối chứng trực tuyến cho S1--S3}
Bộ đối chứng được chọn theo cơ chế: DLS, RDG, TransProtect, semantic correlation và fake-query insertion. Chúng đại diện cho các hướng dummy/vị trí thay thế phù hợp LBS, chưa bao phủ toàn bộ hướng bảo vệ như mix-zone, phối hợp nhiều người hoặc truy vấn riêng tư bằng mật mã.

\par\small
\begin{tabularx}{\linewidth}{@{}P{2.65cm}Y Y@{}}
\toprule\textbf{Phương pháp}&\textbf{Vai trò so sánh}&\textbf{Điều kiện triển khai}\\\midrule
DLS & Đối chứng nền theo xác suất vị trí, không học sâu. & Học thống kê từ train; bản trên mạng đường là thích nghi từ miền gốc.\\
RDG & Đo đóng góp của tương quan chuyển tiếp và chống Viterbi. & Cần thống kê chuyển tiếp; chưa chấm trên tập mới.\\
TransProtect & Một điểm thay thế theo ngữ cảnh đường/chuyển động. & Có lõi và kiến trúc học tùy chọn; benchmark cũ dùng Markov proxy, chưa tương đương mô hình gốc.\\
Semantic correlation & Dummy theo thời gian và ngữ nghĩa. & Adapter hiện dùng loại đường/tần suất thay POI/LSTM gốc; cần hoàn thiện dữ liệu và mô hình.\\
Fake-query insertion & Can thiệp lịch công bố bằng bản tin giả. & Hỗ trợ thời điểm chèn; không lộ nhãn bản tin giả; tính toàn bộ chi phí.\\\bottomrule
\end{tabularx}
\normalsize
\textbf{Nhánh bổ sung.} AnotherMe [20] dùng quỹ đạo ảo; adapter hiện xử lý cả hành trình nên đặt ở nhánh ngoại tuyến. Bổ sung không bảo vệ và dummy ngẫu nhiên làm mốc kiểm tra. Các bản thích nghi được phân biệt với tái lập phương pháp gốc.

\subsection{Metrics gốc và phần được kế thừa}
\par\small
\begin{tabularx}{\linewidth}{@{}P{2.55cm}Y P{3.0cm}P{2.55cm}@{}}
\toprule\textbf{Đối chứng}&\textbf{Privacy / chẩn đoán}&\textbf{Utility}&\textbf{Performance}\\\midrule
DLS & Entropy ứng viên $H\up$. & --- & ---\\
RDG & Entropy ứng viên/chuyển tiếp $\up$; tỷ lệ bảo vệ dưới Viterbi $\up$. & --- & ---\\
TransProtect & EIE $\up$: sai số suy luận kỳ vọng. & Sai lệch chi phí hành trình $\Delta c\down$. & ---\\
Semantic correlation & ASR $\up$; DER $\up$ chẩn đoán hiệu quả dummy. & DER không đo chất lượng POI. & Thời gian sinh $\down$.\\
Fake-query insertion & ASR $\up$; số đường khó phân biệt $\up$. & --- & Độ trễ/thời gian sinh, RSS $\down$.\\\bottomrule
\end{tabularx}
\normalsize
Dấu ---: chỉ số gốc chưa được xác minh trong hồ sơ khảo sát.

\textbf{Công thức đại diện.}
\[
H=-\sum_jp_j\log_2p_j,\quad
\mathrm{ASR}=\frac{N_{\mathrm{đạt\ ẩn\ danh}}}{N_{\mathrm{phép\ thử}}},\quad
\Delta c=\sum_p\pi(p)|c(x,p)-c(y,p)|.
\]
Entropy đo bất định của phân bố; ASR đếm số lần đạt điều kiện ẩn danh; $\Delta c$ đo sai lệch chi phí tới cùng đích khi thay vị trí thật $x$ bằng $y$. EIE có ý nghĩa tổng quát $\mathbb E[d(X,\widehat X)]$; cách kỳ vọng và đối thủ phải được xác định khi tái lập.

\key{Kế thừa bốn ý tưởng: bất định, thành công của đối thủ, sai số suy luận và tổn thất dịch vụ; thống nhất phép đo trên tác vụ chung.}
ASR lặp lại giữa hai paper nhưng điều kiện thành công cần đối chiếu. EIE và MAE cùng họ sai số, chưa mặc nhiên cùng định nghĩa. Entropy/DER giữ ở nhánh chẩn đoán; thời gian xử lý được kế thừa với cùng phần cứng và ranh giới đo. Bộ chung không cộng các đại lượng thành một điểm tổng.

\newpage
\subsection{Bộ đo đề xuất theo ba trục}
\textbf{Privacy: suy luận của đối thủ.} Với vị trí thật $x_i$, dự đoán $\widehat x_i$ và khoảng cách thẳng $e_i=d_E(x_i,\widehat x_i)$ theo mét:
\[
\mathrm{MAE}=\frac1n\sum_i e_i\quad(\up),\qquad
\mathrm{Hit}_r=\frac1n\sum_i\mathbf1[e_i\le r]\quad(\down).
\]
MAE đo độ lệch trung bình; Hit100 đo tần suất đối thủ vẫn định vị trong 100 m. Báo thêm median/p90 sai số và Hit50/Hit200. Sai số là giữa \emph{dự đoán của đối thủ} và vị trí thật, không phải khoảng cách dummy--thật. Bán kính 100 m là lựa chọn thực nghiệm.

\par\small
\begin{tabularx}{\linewidth}{@{}P{2.3cm}Y Y@{}}
\toprule\textbf{Kịch bản}&\textbf{Metric privacy}&\textbf{Đơn vị / cách chấm}\\\midrule
S1--S3 & MAE, Hit, phân vị sai số. & Theo thời điểm/cửa sổ; tổng hợp theo chuyến và ca.\\
S4 & Accuracy hoặc F1 liên kết. & Chấm người và thiết bị riêng; giữ tỷ lệ cặp dương/âm.\\
S5--S6 & Accuracy cạnh/đích; sai số tới đích. & Cùng tiền tố và thời điểm dự đoán; không dùng tương lai.\\
S7 & Macro-F1 ý định. & Trung bình đều qua các lớp; tách lộ trực tiếp và suy luận.\\
S8 & Chênh lệch MAE/Hit. & Cùng phép thử, bật/tắt dữ liệu đồng hành.\\
S9--S10 & MAE/Hit điểm đầu/cuối. & Một điểm mục tiêu mỗi chuyến, theo phần bị che.\\\bottomrule
\end{tabularx}
\normalsize
Với nhiệm vụ phân loại, $\mathrm{Accuracy}=N_{\mathrm{đúng}}/N$,
$F1=2TP/(2TP+FP+FN)$ và $\mathrm{MacroF1}=C^{-1}\sum_cF1_c$.
Các điểm của đối thủ càng cao thì riêng tư càng kém. Với tỷ lệ thành công $S$, báo thêm $\Delta S=S(\text{công bố + phụ trợ})-S(\text{chỉ phụ trợ})$ để tách phần lộ thêm do công bố.

\textbf{Utility: chất lượng truy hồi POI.} Gọi $R_i^\star$ là tối đa 5 POI tham chiếu tại vị trí thật, $\widehat R_i$ là tối đa 5 POI trả sau gộp, loại trùng và lọc:
\[
\mathrm{Recall@5}(i)=\frac{|R_i^\star\cap\widehat R_i|}{|R_i^\star|}\quad(\up),\qquad
C_i=\mathbf1[|\widehat R_i|=|R_i^\star|].
\]
Tỷ lệ trả đủ là trung bình $C_i$ trên truy vấn có đáp án. Cần năm POI nhưng giữ được bốn thì Recall bằng 80\%; trả đủ năm chưa có nghĩa đúng năm. Có đáp án nhưng trả rỗng nhận Recall bằng 0; không có đáp án ghi không áp dụng.
\[
\Delta D_i=\frac{\sum_{p\in\widehat R_i}d_G(Q(x_i),Q(p))}{|\widehat R_i|}
-\frac{\sum_{p\in R_i^\star}d_G(Q(x_i),Q(p))}{|R_i^\star|}\quad(\down).
\]
$Q$ ánh xạ tới mạng đường; $d_G$ là khoảng cách ngắn nhất theo chiều đường. Chỉ tính $\Delta D$ trên truy vấn trả đủ và báo mẫu số. Đây là chênh lệch khoảng cách trung bình tới hai tập POI, khác $\Delta c$ của TransProtect.

\textbf{Performance/cost: tài nguyên cần sử dụng.}
\[
C_{\mathrm{msg}}=\frac{N_{\mathrm{bản\ tin}}}{N_{\mathrm{truy\ vấn\ thật}}},\qquad
C_{\mathrm{bytes}}=\frac{B_{\mathrm{yêu\ cầu}}+B_{\mathrm{phản\ hồi}}}{N_{\mathrm{truy\ vấn\ thật}}}.
\]
Đếm riêng truy vấn loại--tọa độ, số tọa độ công bố/phân biệt và số mục phản hồi. Báo thời gian sinh và hoàn tất dịch vụ p50/p95 trên cùng phần cứng; tách học, chuẩn bị, trực tuyến và mạng. Pin/bộ nhớ là phép đo mở rộng theo thiết bị.

\newpage
\subsection{Đánh đổi privacy--utility--cost và quy tắc lựa chọn}
Ba trục đo các yêu cầu đồng thời. Tăng mức bảo vệ có thể giảm utility, tăng chi phí hoặc cả hai; hướng thay đổi phụ thuộc cơ chế và cần được xác nhận thực nghiệm.
\begin{center}
\begin{tikzpicture}[x=1cm,y=1cm,>=Latex,font=\small]
\draw[->,thick] (0,0)--(11.8,0) node[below left]{Hit100 (\%) $\down$};
\draw[->,thick] (0,0)--(0,4.2) node[above right]{Recall@5 (\%) $\up$};
\draw[dashed,gray] (0,2.1)--(11.5,2.1);
\node[anchor=west,fill=white] at (6.4,2.1) {Ngưỡng chất lượng 90\%};
\draw[->,teal,line width=1.3pt] (4.5,2.8)--(1.0,3.6);
\node[align=left,anchor=west,text=teal] at (4.7,3.15) {Hướng mong muốn:\\Hit thấp, Recall cao};
\draw[teal,fill=pale] (7.4,0.95) circle (0.13);
\draw[teal,fill=pale] (8.1,0.95) circle (0.25);
\draw[teal,fill=pale] (9.05,0.95) circle (0.38);
\node[anchor=west] at (5.7,0.28) {Kích thước điểm = chi phí};
\node[anchor=west] at (0.1,-0.65) {Mỗi điểm: một phương pháp và một cấu hình; màu: phương pháp.};
\end{tikzpicture}
\end{center}
\textbf{Hình 2.} Thiết kế biểu đồ đánh đổi cho S1--S3; đây là sơ đồ cách đọc, chưa biểu diễn kết quả thực nghiệm. Khi chấm benchmark, điểm phía trên trái và kích thước nhỏ là cấu hình mong muốn. Mỗi biểu đồ dùng một đại lượng chi phí cụ thể, chẳng hạn byte/truy vấn thật.

\begin{tabularx}{\linewidth}{@{}P{4.0cm}Y@{}}
\toprule\textbf{Thay đổi cấu hình}&\textbf{Quan hệ cần kiểm tra}\\\midrule
Làm mờ mạnh hơn, giữ ngân sách & Hit có thể giảm nhưng Recall cũng có thể giảm.\\
Tăng dummy / truy vấn chèn & Có thể giảm Hit hoặc giữ Recall; số truy vấn và phản hồi tăng.\\
Giảm tần suất công bố & Giảm chi phí truyền; dịch vụ có thể kém cập nhật, hiệu quả riêng tư phụ thuộc tương quan.\\\bottomrule
\end{tabularx}

\textbf{Quy tắc chọn trên tập validation.} Với cấu hình $\theta$, ca con $s$ và bộ đối thủ $\mathcal A_s$:
\[
\min_s\overline{\mathrm{Recall@5}}_{s,\theta}\ge0.90,
\qquad
\theta^*=\arg\min_{\theta\ \mathrm{đạt}}\frac1{|\mathcal S|}\sum_{s\in\mathcal S}
\max_{a\in\mathcal A_s}\mathrm{Hit100}_{a,s,\theta}.
\]
Trước hết từng ca phải đạt chất lượng; sau đó chọn cấu hình có Hit100 trung bình qua ca thấp nhất trước đối thủ mạnh nhất đã thử. Nếu không cấu hình nào đạt, kết luận chưa có cấu hình khả thi trong lưới thử. Ngân sách chi phí hoặc các mức ngân sách so sánh cần được chốt cùng protocol.

\textbf{Điều kiện so sánh.} Cùng dữ liệu, POI, phần quan sát và thông tin phụ trợ; đối thủ được học/chọn riêng phù hợp từng cơ chế rồi cố định khi kiểm thử. Tổng hợp theo ca/chuyến để chuyến dài không chi phối. Phương pháp gửi thật cùng dummy có thể giữ utility cao, nhưng phải tính mọi truy vấn và phản hồi phát sinh.

\key{Ba trục dùng chung; metric privacy thay theo thông tin cần giấu. Không tạo một điểm riêng tư tổng cho S1--S10.}
$K$ là số tọa độ công bố khi áp dụng; $k=5$ là số POI ứng dụng cần; $L$ là số POI máy chủ trả cho mỗi truy vấn. $K$ ứng viên nội bộ của TransProtect không phải số tọa độ nó công bố. Tăng $L$ vẫn chấm Recall@5 và ghi phần phản hồi tăng thêm.

\newpage
\section{Kế hoạch thực hiện và tiến độ}
\subsection{Các bước tiếp theo}
\begin{center}
\begin{tikzpicture}[box/.style={draw=teal,fill=pale,rounded corners=2pt,text width=6.9cm,minimum height=1.15cm,align=left,font=\small},>=Latex]
\node[box] (a) {\textbf{1. Chốt dữ liệu và threat model}\\Nhãn, quyền quan sát, mẫu độc lập, độ phủ và cách chia tập.};
\node[box,below=0.35cm of a] (b) {\textbf{2. Khóa metrics và protocol S1--S3}\\Đối thủ, mẫu số, tổng hợp, ngưỡng utility và ngân sách.};
\node[box,below=0.35cm of b] (c) {\textbf{3. Hoàn thiện đối chứng và phát triển phương pháp}\\Kiểm tra đầu vào nhân quả; tối ưu và ablation trên tập phát triển.};
\node[box,below=0.35cm of c] (d) {\textbf{4. Xác nhận S1--S3, rồi mở rộng S4--S10}\\Kiểm thử độc lập; thêm cơ chế và đối chứng theo mục tiêu.};
\draw[->] (a)--(b);\draw[->] (b)--(c);\draw[->] (c)--(d);
\end{tikzpicture}
\end{center}
\textbf{Hình 3.} Đặc tả đủ S1--S10 từ đầu; triển khai và xác nhận theo giai đoạn. Mỗi lần mở rộng phải có đầu tấn công và tiêu chí dịch vụ tương ứng.

\subsection{Hướng phương pháp và thực nghiệm hiện có}
BR-Dummy là hướng nghiên cứu truy vấn giả có xét chuyển động trên mạng đường và độ phủ POI của cả tập truy vấn. Thiết bị nhận, gộp và lọc kết quả để phục vụ nhu cầu thật. Giai đoạn S1--S3 kiểm tra giả thuyết rằng lựa chọn tập truy vấn theo ích lợi chung có thể cải thiện đánh đổi riêng tư--chất lượng--chi phí.

Đã có pipeline mô phỏng, sinh bản ghi kịch bản, truy vấn POI và đánh giá một số đối thủ. Thực nghiệm thăm dò S1--S3 so sánh các biến thể nội bộ để kiểm tra bộ đo và nhận diện đánh đổi. Chưa cấu hình nào đạt đầy đủ quy tắc chọn của vòng đã chạy; chưa có đối chứng ngoài trên tập xác nhận mới. Các đầu suy luận S4--S10 chưa có kết quả xác nhận.

Vòng hiện tại đã có MAE, Hit100, Recall@5, tỷ lệ trả đủ, số truy vấn, số mục phản hồi và byte mã POI. Cần bổ sung phân bố sai số, lưu lượng đầy đủ và thời gian xử lý theo protocol. Dữ liệu đã dùng để điều chỉnh thiết kế tiếp tục thuộc tập phát triển; kết luận cuối dùng dữ liệu xác nhận độc lập.

\subsection{Định hướng triển khai}
\begin{enumerate}
\item \textbf{Phạm vi:} S1--S10 là khung bài toán; S1--S3 là mục tiêu phát triển và benchmark đầu tiên.
\item \textbf{Đối chứng:} năm phương pháp trực tuyến đã chọn; AnotherMe ở nhánh bổ sung; mở rộng survey theo mục tiêu bảo vệ.
\item \textbf{Đánh giá:} Recall@5 từng ca tối thiểu 90\%, Hit100 là tiêu chí chọn privacy, MAE và các phân vị hỗ trợ; chốt các mức ngân sách chi phí.
\item \textbf{Mở rộng:} ưu tiên theo rủi ro còn lại và bằng chứng S1--S3, bổ sung cơ chế cho định danh, nội dung, đồng hành và hai đầu hành trình.
\end{enumerate}
\key{Ưu tiên trước mắt: khóa dataset và giao thức S1--S3, hoàn thiện đối chứng, rồi tối ưu và xác nhận phương pháp.}

\newpage
\section*{Tài liệu tham khảo}
\small\RaggedRight
\begin{enumerate}[label={[\arabic*]},leftmargin=2em,itemsep=6pt]
\item Niu et al. (2014). \emph{Achieving k-Anonymity in Privacy-Aware Location-Based Services}. INFOCOM. \href{https://doi.org/10.1109/INFOCOM.2014.6848002}{DOI}.
\item Andrés et al. (2013). \emph{Geo-Indistinguishability: Differential Privacy for Location-Based Systems}. CCS. \href{https://arxiv.org/abs/1212.1984}{arXiv:1212.1984}.
\item Sun et al. (2017). \emph{ASA: Against Statistical Attacks for Privacy-Aware Users in Location Based Service}. FGCS 70. \href{https://doi.org/10.1016/j.future.2016.06.017}{DOI}.
\item \emph{A trajectory privacy protection method using cached candidate result sets} (2024). JPDC 193, 104965. \href{https://doi.org/10.1016/j.jpdc.2024.104965}{DOI}.
\item Shaham et al. (2021). \emph{Privacy Preservation in Location-Based Services: A Novel Metric and Attack Model}. IEEE TMC 20(10). \href{https://doi.org/10.1109/TMC.2020.2993599}{DOI}.
\item Yadav et al. (2024). \emph{Protecting Vehicle Location Privacy with Contextually-Driven Synthetic Location Generation}. SIGSPATIAL. \href{https://doi.org/10.1145/3678717.3691211}{DOI}.
\item Liu, Peng, Zhou (2026). \emph{A Dummy-Based Location Privacy Protection Scheme with Semantic Correlation of Moving Paths}. JKSUCIS 38:478. \href{https://doi.org/10.1007/s44443-026-00899-w}{DOI}.
\item Liu, Hu, Zhou (2026). \emph{Location Privacy Protection in Continuous LBSs: Enhancing Anonymity via Fake Queries}. JKSUCIS 38:53. \href{https://doi.org/10.1007/s44443-025-00438-z}{DOI}.
\item Beresford, Stajano (2004). \emph{Mix Zones: User Privacy in Location-aware Services}. PerCom Workshops. \href{https://doi.org/10.1109/PERCOMW.2004.1276918}{DOI}.
\item Wang et al. (2025). \emph{CPCROK: A Communication-Efficient and Privacy-Preserving Scheme for Low-Density Vehicular Ad Hoc Networks}. Future Internet 17(4):165. \href{https://doi.org/10.3390/fi17040165}{DOI}.
\item Theodorakopoulos et al. (2014). \emph{Prolonging the Hide-and-Seek Game: Optimal Trajectory Privacy for Location-Based Services}. WPES. \href{https://arxiv.org/abs/1409.1716}{arXiv:1409.1716}.
\item Qiu et al. (2023). \emph{Novel trajectory privacy protection method against prediction attacks}. ESWA 213, 118870. \href{https://doi.org/10.1016/j.eswa.2022.118870}{DOI}.
\item Xue et al. (2017). \emph{A Novel Destination Prediction Attack and Corresponding Location Privacy Protection Method in Geo-social Networks}. IJDSN 13(1). \href{https://doi.org/10.1177/1550147716685421}{DOI}.
\item Hakeem et al. \emph{DFPS: A Distributed Mobile System for Free Parking Assignment}. IEEE TMC, công bố trực tuyến 2021. \href{https://doi.org/10.1109/TMC.2021.3080222}{DOI}.
\item Wu et al. (2021). \emph{Constructing Dummy Query Sequences to Protect Location Privacy and Query Privacy in Location-Based Services}. WWW 24. \href{https://doi.org/10.1007/s11280-020-00830-x}{DOI}.
\item \emph{A Location Privacy and Query Privacy Joint Protection Scheme for POI Query in Vehicular Networks} (2024). Journal of Electronics \& Information Technology. \href{https://doi.org/10.11999/JEIT221599}{DOI}.
\item Olteanu et al. (2017). \emph{Quantifying Interdependent Privacy Risks with Location Data}. IEEE TMC 16(3). \href{https://doi.org/10.1109/TMC.2016.2561281}{DOI}.
\item \emph{DP-FETC: A differentially private trajectory publishing method based on feature extraction and trajectory correlation} (2025). ERA 33(11). \href{https://doi.org/10.3934/era.2025293}{DOI}.
\item Dhondt et al. (2022). \emph{A Run a Day Won't Keep the Hacker Away: Inference Attacks on Endpoint Privacy Zones in Fitness Tracking Social Networks}. CCS. \href{https://doi.org/10.1145/3548606.3560616}{DOI}.
\item Li et al. (2024). \emph{AnotherMe: A Location Privacy Protection System Based on Online Virtual Trajectory Generation}. IEEE TDSC 21(4). \href{https://doi.org/10.1109/TDSC.2023.3314200}{DOI}.
\end{enumerate}
\textbf{Nguồn dữ liệu và thực nghiệm.} Đặc tả: \path{thesis/scenario_dataset_spec.tex}; metrics: \path{thesis/metrics_consolidation.tex}; protocol: \path{thesis/notes/fresh_switching_protocol.md}; kết quả: \path{artifacts/benchmarks/fresh_switching/readout.json} và \path{selection.json}.
\end{document}
